\chapter{Trabajos futuros}
\label{chap:trabajos-futuros}

El sistema desarrollado cumple los objetivos planteados, sin embargo está abierto a distintas líneas de investigación.
Este capítulo recoge las más relevantes, desde optimizaciones de bajo nivel hasta nuevas capacidades que ampliarían el alcance del proyecto.

\section{Construcción del volumen en la GPU}

La construcción del \acrshort{svo} se ejecuta actualmente en la \acrshort{cpu}, paralelizada mediante el Job System y el compilador Burst.
Dado que la rasterización es un proceso completamente paralelizable, se plantea trasladar la voxelización y la construcción de la jerarquía a la \acrshort{gpu} mediante \textit{compute shaders}~\cite{UnityComputeShaders}.

Esta vía se descartó en el desarrollo por varios motivos.
En primer lugar, la programación de la \acrshort{gpu} no está del todo integrada con la capa de abstracción que ofrece la \acrshort{api} de Unity, de manera que un cambio en estas interfaces podría dejar el sistema obsoleto o comprometer la portabilidad entre distintas plataformas.
En segundo lugar, se trata de una alternativa de riesgo elevado y coste de desarrollo considerable.
Esto se debe a que el código que se ejecuta en la tarjeta gráfica es muy difícil de depurar y además obligaría a mantener un espejo de la representación del volumen en la memoria de la tarjeta gráfica y a sincronizarlo con el de la \acrshort{cpu}, lo que puede dar lugar a errores de consistencia entre ambas representaciones.
Finalmente, a pesar de existir algoritmos muy rápidos con este fin~\cite{schwarz2010fast}, la construcción habitualmente se realiza fuera del tiempo de edición y, tal como se comentó en la Sección~\ref{sec:mediciones}, esta suele durar menos de cinco segundos usando el máximo de precisión, por lo que en ningún momento supone el cuello de botella del sistema.

En conclusión, se trata de una optimización interesante pero de prioridad baja.
Esta solo cobraría sentido si el tiempo de construcción llegase a ser crítico en escenas de gran tamaño o si las herramientas de cálculo en la \acrshort{gpu} del motor alcanzasen una mayor estabilidad.

\section{Búsqueda y evasión mediante aprendizaje por refuerzo}

El cálculo de rutas se apoya en una heurística diseñada a mano, ya sea la distancia euclídea o el conteo de nodos, sobre una búsqueda \gls{astar} determinista, mientras que la evasión recae en el algoritmo \acrshort{orca}.
Ambos son enfoques reactivos y deterministas.
Una línea de investigación que se propone consiste en incorporar técnicas de aprendizaje automático, en concreto de \textbf{aprendizaje por refuerzo}, para mejorar la búsqueda de caminos y la evasión entre agentes.

Esta idea parte de investigaciones previas de aplicar este conjunto de técnicas al control de enjambres heterogéneos de aeronaves pilotadas remotamente, coloquialmente conocidos como drones~\cite{puente_castro_2023}.
Implementarlas en el sistema, podría no solo mejorar la eficiencia del mismo al evitar ejecutar las partes con mayor carga algorítmica, sino que tambíen proporcionaría resultados más naturales y creíbles.
Esto se debe a que los agentes considerarán y evaluarán todo tipo de variables, incluyendo las que no se hayan tenido en cuenta durante el desarrollo de este proyecto.

\section{Otras líneas de extensión}

Además de las dos grandes líneas anteriores, el proyecto deja abiertas otras posibles mejoras:

\begin{description}
    \item \textbf{Geometría dinámica y reconstrucción incremental:} reconstruir únicamente las regiones del \acrshort{svo} afectadas por un cambio en la geometría estática, como una puerta que se abre o una estructura que se destruye, superando la limitación actual de un volumen totalmente precalculado y orientado a entornos estáticos.
    \item \textbf{Búsqueda jerárquica:} para aliviar el coste del cálculo de rutas en escenas densas, explorar esquemas de búsqueda jerárquica o de cualquier ángulo (\textit{any-angle})~\cite{harabor2013optimal}, así como la precomputación de la conectividad entre regiones, con el fin de reducir el número de nodos explorados.
    \item \textbf{Enlaces entre volúmenes:} en escenarios de gran tamaño, resulta ineficiente mantener un único volumen gigante cargado en memoria.
    Una mejora sería permitir a los agentes desplazarse de forma fluida entre diferentes volúmenes independientes lo que implicaría desarrollar un sistema de enlaces que conecten las fronteras de estos.
\end{description}
